% ===== §6 =====
\section{The Third Impossibility: Incommensurability}
\label{p25:sec:incommensurability}

\noindent
Grant once more, for the sake of pressing on, that the objective could be written and that some standpoint could fix it. The third condition of the anatomy then requires that the values the objective assigns lie on a common scale, so that any two candidates compare and $\argmax$ has something to select. This section denies that condition through the incommensurability of certain goods: the thesis, developed most rigorously by Raz, that there are goods between which no common measure holds, so that the demand to rank them on one scale is not a hard task but a mistaken one, and that within intimate relation the attempt to price a good can be, in the strict sense, a betrayal of it rather than a misjudgment. As with the previous two impossibilities, the claim must first be distinguished from a weaker neighbor with which the frame would gladly confuse it, and only then pressed to its sharp form.

\subsection*{Incommensurability as a relation, not a difficulty}

The weaker neighbor is the thought that some comparisons are merely hard, that weighing a career against a friendship is difficult and uncertain and beyond our present precision. This the optimizing frame absorbs without strain: hard comparisons are comparisons still, and difficulty is only a call for better estimates. Incommensurability in the strict sense is a different matter, a positive structural relation and not a confession of ignorance. Two goods are incommensurable when there is no common measure in terms of which their values are precisely comparable, when it is neither true that the first is better than the second, nor that the second is better than the first, nor that they are exactly equal \parencite{raz1986morality,chang1997incommensurability}. The three familiar relations of a common scale, better, worse, equal, are jointly exhausted only where a scale underlies them; where none does, two goods may stand outside all three, related in a way the scale cannot represent, and this standing-outside is not our failure to have measured but the goods' failure to share a measure. The most exact mark of the relation is due to Raz and deserves to be given precisely, because it converts a vague intuition into a test. Suppose two options were merely \emph{equal} on some common scale. Then improving one of them even slightly, sweetening it by a little, would have to make it determinately better than the other, since on a scale a small addition to an equal quantity yields a larger quantity. Now take a career and a friendship one is inclined to call equally choiceworthy, and sweeten the career by a modest raise in salary. If the two were equal on a scale, the raised career is now determinately the better choice. But it is not: the career with the raise is still neither better than the friendship, nor worse, nor exactly equal, just as before the raise. The tie does not break under the small improvement, and a tie that does not break under improvement was never a tie on a scale at all. The small-improvement argument shows, not that the comparison is hard, but that there is no scale beneath it to be hard about.

\subsection*{Pricing as betrayal}

Within intimate relation this structural relation takes a sharper form, in which the attempt at commensuration is not only a mistake about the goods but a wrong done to them, and here the paper leans on the line running from Raz through Anderson to the recent moral critique of markets \parencite{anderson1993value,sandel2012money}. The sharpening turns on a fact about certain goods: that they are constituted, in part, by the stances one takes toward them, so that to take the wrong stance is not to misdescribe the good but to alter it. Among the goods so constituted are those whose being includes a refusal to price them. For such a good, to price it is not to measure it, even badly; it is to change it into something else, and in the changing to betray it. This is why the refusal to name a sum is not squeamishness. To ask what payment would compensate the ending of a friendship is not to arrive, with effort, at a large figure; it is to have misunderstood what a friendship is, since a bond one holds oneself ready to release for a price is, by that readiness, already not the bond one took oneself to be pricing but a lesser thing wearing its name. The pricing does not report a worth and get the figure wrong. It produces, in the act of pricing, a diminished thing, and reports a fact about that.

\begin{casebox}{the offer}
A man is asked, by someone who means the question seriously, what sum would induce him to give up his oldest friendship, permanently and with no further contact. Two things are worth noticing in his response. The first is that he does not experience the question as difficult, as though the number were very large and hard to fix; he experiences it as a category mistake, as one would the question how much the number seven weighs. The second is more telling. If he were to take the question in the spirit asked, to begin in earnest to name a figure and to negotiate it upward, something would already have happened to the friendship, before any sum was agreed and whether or not one ever was. To have entertained the offer in earnest is to have constituted oneself as someone whose friendships are for sale at a price, and a friendship held by such a person is not the good the offer was meant to price. The betrayal is not in the accepting, which never comes; it is in the pricing, which changes the thing priced in the moment it treats it as priceable. This is why the fitting answer is not a very large number but a refusal of the frame, and why the refusal is not a failure to compute but the only response that leaves the good intact.
\end{casebox}

\begin{claim}[Incommensurability]
Two goods are \emph{incommensurable} when they stand in none of the three relations a common scale affords, being neither better, nor worse, nor exactly equal, as the failure of their tie to break under small improvement reveals. Certain goods at stake in intimate relation are incommensurable in this sense; and some are further constituted, in part, by the refusal to price them, so that to place them on a common scale is not to mismeasure them but to transform and thereby betray them. The scalar condition of the anatomy therefore fails: the values the objective assigns do not all lie on a common scale, and $\argmax$, which requires such a scale to have anything to select, has nothing over which to operate. The error is one of category and not of computation; the goods are not the kind of thing of which there is a sum.
\end{claim}

\subsection*{The relation to the first two impossibilities}

The three impossibilities so far compound rather than repeat one another, and it is worth marking the joints at which each cuts, since an objector may suspect that they are a single complaint in three costumes. The first concerned \emph{representation}: whether a value can enter the objective as a term at all, and it found that the drive cannot, being extinguished by the attainment a monotone magnitude would track. The second concerned \emph{fixation}: whether the objective, granting it written, can be given from an exterior standpoint, and it found that relational value admits none, being generated within the relation that bears it. This third concerns \emph{comparison}: whether the values the objective assigns, granting them written and fixed, lie on the common scale that maximization requires, and it finds that incommensurable goods do not. The three are independent, and their independence is the strength of the case, for an objector who somehow answered one would face the others untouched: were the drive representable after all, its term would still have to be commensurated with the rest, and incommensurability would block it; were an exterior standpoint somehow found to fix the objective, the goods it fixed would still fail to share a scale. The descent is real, from whether the problem can be written, through whether it can be fixed, to whether what it fixes can be compared; and beneath all three lies a fourth question the descent has been circling from the start, not yet touched, the question of the one for whom the problem is a problem at all, the optimizer whose stable preferences the whole construction has quietly presupposed and whom the next section dissolves.
